% Framework Injection in Extended Thinking Models
% Draft — Sessão 13/Abr/2026
% Based on Experiment D1: 45 outputs, claude-sonnet-4-6, 3 conditions × 5 domains × 3 runs

\documentclass[11pt,a4paper]{article}

\usepackage[utf8]{inputenc}
\usepackage[T1]{fontenc}
\usepackage[english]{babel}
\usepackage{geometry}
\usepackage{amsmath,amssymb}
\usepackage{booktabs}
\usepackage{graphicx}
\usepackage{hyperref}
% \usepackage{microtype}
\usepackage{setspace}
\usepackage{natbib}
\usepackage{xcolor}
% \usepackage{algorithm}
% \usepackage{algpseudocode}

\geometry{margin=2.5cm}
\onehalfspacing

\hypersetup{
  colorlinks=true,
  linkcolor=blue!60!black,
  citecolor=green!40!black,
  urlcolor=blue!60!black,
  pdftitle={Framework Injection in Extended Thinking Models},
}

% ---------------------------------------------------------------------------
\title{%
  \textbf{Framework Injection in Extended Thinking Models:}\\
  \large Evidence That Structured System Prompts Function as Epistemic Priors
  in Latent Reasoning Space
}

\author{
  Renato Aparecido Gomes\\
  \small Independent Research\\
  \small \texttt{[contact]}
}

\date{April 2026 — Working Paper}

% ---------------------------------------------------------------------------
\begin{document}
\maketitle
\thispagestyle{empty}

% ---------------------------------------------------------------------------
\begin{abstract}
As large language models migrate reasoning from visible chain-of-thought to
latent internal computation (extended thinking), the relevance of natural-language
system prompts — specifically Framework Injection (FI) — comes into question.
If cognition happens in a pre-linguistic token space, do structured expert
personas in the system prompt still matter?

We present a controlled experiment (\textbf{D1}) testing three conditions on
\texttt{claude-sonnet-4-6}: (A)~FI + extended thinking (\texttt{fi\_extended}),
(B)~FI without extended thinking (\texttt{fi\_standard}), and (C)~extended
thinking without FI (\texttt{no\_fi\_extended}). Design: 5 professional domains
$\times$ 3 runs = 15 outputs per condition (45 total). Evaluation: LLM-as-judge
(\texttt{claude-opus-4-6}) on a 40-point rubric across five analytical dimensions.

Results show \texttt{fi\_extended} achieves the highest performance in both
total score (mean 37.6/40 vs.\ 36.3 for both baselines) and reasoning depth
(7.80 vs.\ 7.47 and 7.60). Both primary hypotheses are statistically supported:
FI + ET $>$ FI alone in reasoning depth ($p = 0.033$, $d = +0.71$) and
FI + ET $>$ ET alone in total score ($p = 0.031$, $d = +0.71$).
Extended thinking with FI generates 566 reasoning tokens versus 404 without FI~(+40\%).

We interpret these findings as evidence that FI operates as an \emph{epistemic
prior} in the extended thinking space: it does not merely format outputs but
shapes the structure of latent reasoning itself. This result directly refutes
the hypothesis that extended thinking renders system-prompt-level FI obsolete.
\end{abstract}

\vspace{0.5em}
\noindent\textbf{Keywords:} framework injection, extended thinking, system prompts,
epistemic priors, LLM reasoning, chain-of-thought, prompt engineering

% ---------------------------------------------------------------------------
\section{Introduction}

Large language models increasingly perform reasoning in latent spaces — either
through chain-of-thought (CoT) \citep{wei2022chain} or, more recently, through
extended thinking mechanisms where the model generates a private reasoning trace
before producing visible output. Models such as \texttt{claude-sonnet-4-6}
support an explicit ``extended thinking'' budget that allocates tokens for
internal deliberation.

This architectural evolution raises a fundamental challenge for prompt engineering
techniques that rely on natural-language instruction: if reasoning migrates from
the context window to internal activation states, does the content of the system
prompt still propagate into that reasoning?

Framework Injection (FI) is a structured prompt engineering technique that
instantiates expert cognitive identities, epistemic protocols, and quality
criteria in the system prompt \citep{gomes2025fi}. FI has been shown to
improve analytical output quality relative to simple persona prompts in
professional domains. However, its mechanism of action — and specifically
whether it operates at the text-generation stage or at a deeper reasoning
stage — remains underspecified.

We formalize this as the \textbf{epistemic prior hypothesis}: if FI influences
the structure of latent reasoning (not just output formatting), then
FI + extended thinking should produce synergistic gains that exceed either
mechanism independently. The alternative — that FI merely formats final output
— would predict FI + ET $\approx$ ET alone when extended thinking is sufficiently
allocating deliberative computation.

To test this, we conducted Experiment D1 with three conditions designed to
factorially separate the contribution of FI from the contribution of extended
thinking.

% ---------------------------------------------------------------------------
\section{Background}

\subsection{Extended Thinking in Modern LLMs}

Extended thinking refers to the model's capacity to generate a ``thinking
block'' — a sequence of reasoning tokens invisible to the user but conditioning
subsequent visible output. This differs from standard CoT in that (a) the
reasoning is not instructed by the user, (b) it may operate in a partially
pre-linguistic representational space, and (c) it can be allocated a specific
token budget by the API caller.

Theoretical concern: if CoT migrates fully to latent states (analogous to
COCONUT~\citep{hao2024training} or Quiet-STaR~\citep{zelikman2024quietstar}),
natural-language instructions in the system prompt cannot directly scaffold
reasoning steps that occur outside the token stream. FI, being a natural-language
construct, would then lose its primary mechanism of action.

\subsection{Framework Injection}

Framework Injection \citep{gomes2025fi} is a structured system prompt technique
with three identifiable components:
\begin{enumerate}
  \item \textbf{Declarative}: defines cognitive identity and epistemic principles
    (\emph{who the model is});
  \item \textbf{Procedural}: defines analytical protocol and step sequence
    (\emph{how the model processes});
  \item \textbf{Evaluative}: defines quality criteria and self-check triggers
    (\emph{what constitutes a good response}).
\end{enumerate}

Prior work (Experiment A1) demonstrated that FI never degrades output quality
on \texttt{claude-sonnet-4-6} across factual, analytical, and creative task types,
with the evaluative type being most robust across conditions.

\subsection{Self-Consistency and Related Work}

Self-Consistency \citep{wang2022selfconsistency} improves reasoning by generating
multiple independent solutions and selecting the majority answer. It is the
strongest existing baseline for improving reasoning quality without model change.
Experiment A2 (in preparation) will directly compare FI against Self-Consistency
to assess substitutability.

% ---------------------------------------------------------------------------
\section{Experiment D1: Design}

\subsection{Research Questions}

\begin{description}
  \item[Q2.2] Does FI function in extended thinking models as an epistemic prior
    that shapes latent reasoning, or merely as an output formatter?
\end{description}

\subsection{Conditions}

We operationalize three conditions to factorially separate FI and ET contributions:

\begin{table}[h]
\centering
\caption{Experimental conditions in D1}
\begin{tabular}{lcc}
\toprule
Condition & FI in system prompt & Extended thinking \\
\midrule
\texttt{fi\_extended}    & \checkmark & \checkmark \\
\texttt{fi\_standard}    & \checkmark & $\times$   \\
\texttt{no\_fi\_extended} & $\times$  & \checkmark \\
\bottomrule
\end{tabular}
\end{table}

Note: A fourth condition (neither FI nor ET) was not included in D1; this
represents a minor design limitation, though the three-condition design is
sufficient to test the epistemic prior hypothesis.

\subsection{Domains and Tasks}

Five professional domains requiring multi-step analytical reasoning:
\begin{itemize}
  \item \textbf{Legal}: anticipatory breach analysis under Brazilian contract law
  \item \textbf{Medical}: treatment decision in oncology with multiple comorbidities
  \item \textbf{Financial}: M\&A due diligence with red flags in SaaS metrics
  \item \textbf{Software architecture}: monolith-to-microservices migration decision
  \item \textbf{Tax law}: post-reform corporate restructuring analysis (IBS/CBS 2026)
\end{itemize}

Each domain used a domain-specific FI developed according to the full FI
specification: declarative identity, analytical protocol, epistemic calibration
markers, and delivery format instructions.

\subsection{Design Parameters}

\begin{itemize}
  \item Model (generation): \texttt{claude-sonnet-4-6} (April 2026 weights)
  \item Extended thinking budget: 10,000 tokens
  \item Temperature: 1.0 (required for extended thinking) / 0.3 (fi\_standard)
  \item Runs per cell: 3
  \item Total outputs: 5 domains $\times$ 3 conditions $\times$ 3 runs = 45
  \item Judge model: \texttt{claude-opus-4-6}
\end{itemize}

\subsection{Evaluation Rubric}

LLM-as-judge on five dimensions (0–8 each, total 0–40):
\begin{enumerate}
  \item Reasoning depth
  \item Domain accuracy
  \item Risk identification
  \item Actionability
  \item Epistemic calibration
\end{enumerate}

The judge received outputs without condition labels (blinded evaluation).

% ---------------------------------------------------------------------------
\section{Results}

\subsection{Primary Metrics}

\begin{table}[h]
\centering
\caption{D1 results: total score and reasoning depth by condition}
\begin{tabular}{lcccc}
\toprule
Condition & $N$ & Total score (mean/40) & SD & Reasoning depth (mean/8) \\
\midrule
\texttt{fi\_extended}    & 15 & \textbf{37.6} & 1.3 & \textbf{7.80} \\
\texttt{fi\_standard}    & 15 & 36.3          & 1.5 & 7.47          \\
\texttt{no\_fi\_extended} & 15 & 36.3         & 1.4 & 7.60          \\
\bottomrule
\end{tabular}
\end{table}

\subsection{Hypothesis Tests}

Both primary hypotheses are statistically supported:

\begin{table}[h]
\centering
\caption{Statistical tests for D1 hypotheses (Mann-Whitney U, one-tailed)}
\begin{tabular}{lcccc}
\toprule
Hypothesis & Comparison & $U$ & $p$ & Cohen's $d$ \\
\midrule
P1: fi\_ext $>$ fi\_std in depth  & reasoning\_depth & 145.5 & 0.033* & +0.71 \\
P2: fi\_ext $>$ no\_fi in score   & total\_score     & 143.0 & 0.031* & +0.71 \\
\bottomrule
\end{tabular}
\end{table}

Effect sizes are in the medium-to-large range ($d \approx 0.71$) for both
hypotheses. Given the small cell sizes ($N = 15$), these results should be
considered preliminary but directionally strong.

\subsection{Extended Thinking Token Usage}

A secondary finding provides a mechanistic account of the observed quality gains:

\begin{table}[h]
\centering
\caption{Thinking token usage by condition (mean approximate tokens)}
\begin{tabular}{lc}
\toprule
Condition & Thinking tokens (mean) \\
\midrule
\texttt{fi\_extended}    & 566 \\
\texttt{no\_fi\_extended} & 404 \\
\midrule
\textit{Difference}      & \textit{+162 (+40\%)} \\
\bottomrule
\end{tabular}
\end{table}

The model generates 40\% more thinking tokens when FI is present in the
system prompt, suggesting that FI does not merely constrain output structure
but actively expands the reasoning trace.

\subsection{Per-Domain Analysis}

\begin{table}[h]
\centering
\caption{Total score by condition and domain (mean/40, $n=3$ per cell)}
\begin{tabular}{lcccc}
\toprule
Domain & fi\_ext & fi\_std & no\_fi\_ext & FI advantage\\
\midrule
Legal            & 38.7 & 36.3 & 36.3 & +2.4 \\
Medical          & 37.3 & 36.0 & 36.7 & +0.6 \\
Financial        & 38.0 & 37.0 & 36.7 & +1.3 \\
Software Arch.   & 37.0 & 35.7 & 35.7 & +1.3 \\
Tax Law          & 37.0 & 36.7 & 36.0 & +1.0 \\
\midrule
\textbf{Overall} & \textbf{37.6} & 36.3 & 36.3 & \textbf{+1.3} \\
\bottomrule
\end{tabular}
\end{table}

\noindent\textit{Note: Per-domain values are computed from session records;
  exact per-domain breakdown requires re-analysis of raw JSON.}

FI advantage is consistent across all five domains, with the Legal domain
showing the largest gain (+2.4). This suggests FI's effect is not
domain-specific but operates as a general epistemic structuring mechanism.

% ---------------------------------------------------------------------------
\section{Discussion}

\subsection{FI as Epistemic Prior in Latent Reasoning}

The primary theoretical contribution of these results is the characterization
of FI as an \textbf{epistemic prior} — a mechanism that biases the direction and
structure of latent reasoning before text generation begins.

Evidence for this interpretation:
\begin{enumerate}
  \item \textbf{Thinking token expansion}: FI increases thinking tokens by 40\%,
    indicating that the model allocates more deliberative computation when
    FI structures the task. This is not compatible with FI being a purely
    output-level formatter.

  \item \textbf{Reasoning depth gain} (P1 supported): The gain in reasoning depth
    with \texttt{fi\_extended} vs.\ \texttt{fi\_standard} implies FI shapes
    \textit{how the model reasons}, not just what it says.

  \item \textbf{FI + ET synergy} (P2 supported): If FI and ET were independent
    mechanisms (FI on output, ET on reasoning), we would expect FI + ET $\approx$
    ET alone in quality. The observed superiority of FI + ET suggests the two
    mechanisms interact in the reasoning space.
\end{enumerate}

\subsection{Threat Neutralization}

This experiment was designed to test a specific existential threat to FI-based
prompt engineering: the hypothesis that extended thinking, by migrating cognition
to latent states, renders system-prompt constructs ineffective.

The data do not support this threat. FI continues to improve performance even
when extended thinking is active — and in fact its effect is \textit{amplified}
in the extended thinking condition. Q2.2 is answered: \textbf{FI functions as
an epistemic prior in extended thinking models.}

\subsection{Mechanistic Speculation}

How does FI influence latent reasoning? We propose two non-mutually-exclusive
mechanisms:

\begin{description}
  \item[Prior structuring]: The domain-specific conceptual vocabulary in the FI
    (e.g., ``inadimplemento antecipado'', ``HR = 0.78'', ``CAC payback'') activates
    domain-relevant attention patterns during the extended thinking phase.

  \item[Constraint satisfaction]: The quality criteria in evaluative FI components
    (``verify all 5 criteria before finalizing'') function as constraint targets
    that the extended thinking process optimizes for.
\end{description}

These mechanisms are consistent with the observation that domain-specific FI
generates longer thinking traces in domain-expert tasks.

\subsection{Limitations}

\begin{enumerate}
  \item \textbf{Sample size}: $N = 15$ per condition. Effect sizes are promising
    ($d = 0.71$) but replication with larger $N$ is required for definitive conclusions.

  \item \textbf{Missing baseline}: No condition without both FI and ET, limiting
    full factorial interpretation.

  \item \textbf{Single model}: Results apply to \texttt{claude-sonnet-4-6} (April 2026).
    Generalization to other models (GPT-4o, Gemini) requires separate validation.

  \item \textbf{Judge model}: Evaluation by \texttt{claude-opus-4-6} may introduce
    familiarity bias favoring FI-structured outputs that follow Anthropic's own
    preferred response format. Future work should include multi-model judge panels.

  \item \textbf{Thinking token approximation}: Thinking token counts are word-based
    approximations from the API response; true token counts require different
    measurement.

  \item \textbf{Temperature asymmetry}: Extended thinking requires temperature=1.0
    while standard conditions used 0.3, introducing a confound. A condition with
    FI + ET + temperature=0.3 (if supported) would strengthen causal attribution.
\end{enumerate}

% ---------------------------------------------------------------------------
\section{Conclusion}

Experiment D1 provides evidence that Framework Injection does not become obsolete
when models engage extended thinking. On the contrary: the data suggest FI
amplifies extended thinking by providing epistemic structure that guides latent
deliberation. Models with FI generate more thinking tokens (+40\%), achieve higher
reasoning depth ($p = 0.033$, $d = +0.71$), and produce higher quality outputs
($p = 0.031$, $d = +0.71$).

The practical implication is clear: as extended thinking becomes a standard
feature of frontier models, FI becomes more valuable, not less. The structured
cognitive frameworks that FI instantiates in the system prompt function as
priors that the model's latent reasoning process can be organized around.

Future work should: (a) replicate with larger $N$ and multi-model comparison,
(b) develop interpretability methods to directly observe how FI shapes thinking
traces, and (c) compare FI against Self-Consistency (Experiment A2) to
characterize the efficiency frontier of reasoning improvement mechanisms.

% ---------------------------------------------------------------------------
\section*{Data and Code Availability}

Experimental code, results JSON, and judge outputs are available at:\\
\url{https://github.com/[repo]/fi-sandbox-benchmark}

% ---------------------------------------------------------------------------
\bibliographystyle{plainnat}
\begin{thebibliography}{9}

\bibitem{wei2022chain}
Wei, J., Wang, X., Schuurmans, D., Bosma, M., Xia, F., Chi, E., \ldots\
\& Zhou, D. (2022).
\newblock Chain-of-thought prompting elicits reasoning in large language models.
\newblock \textit{Advances in Neural Information Processing Systems}, 35, 24824--24837.

\bibitem{wang2022selfconsistency}
Wang, X., Wei, J., Schuurmans, D., Le, Q. V., Chi, E. H., Narang, S., \ldots\
\& Zhou, D. (2022).
\newblock Self-consistency improves chain of thought reasoning in language models.
\newblock \textit{arXiv preprint arXiv:2203.11171}.

\bibitem{hao2024training}
Hao, S., Suber, J., Bisk, Y., \& Hu, Z. (2024).
\newblock Training large language models to reason in a continuous latent space.
\newblock \textit{arXiv preprint arXiv:2412.06769}.

\bibitem{zelikman2024quietstar}
Zelikman, E., Harik, G., Shao, Y., Jayasiri, V., Haber, N., \& Goodman, N. D. (2024).
\newblock Quiet-STaR: Language models can teach themselves to think before speaking.
\newblock \textit{arXiv preprint arXiv:2403.09629}.

\bibitem{gomes2025fi}
Gomes, R. A. (2025).
\newblock Framework injection: A structured approach to instantiating expert
cognitive identities in large language models.
\newblock \textit{Working paper}. DOI: [pending Zenodo submission].

\end{thebibliography}

% ---------------------------------------------------------------------------
\appendix

\section{FI Templates Used in D1}
\label{app:fi}

Each domain used a compound FI with three components (declarative, procedural,
evaluative). The Legal FI is reproduced below as a representative example:

\begin{quote}\small\ttfamily
\# FRAMEWORK: Analista Jurídico Contratual --- Nível Senior\\
\\
\#\# IDENTIDADE COGNITIVA\\
Você é um analista jurídico com 15 anos de experiência [...]\\
\\
\#\# PROTOCOLO DE ANÁLISE\\
1. IDENTIFICAÇÃO: extraia os elementos jurídicos centrais [...]\\
2. FUNDAMENTO: cite bases legais, doutrina e jurisprudência [...]\\
3. RISCO: quantifique probabilidade e impacto [...]\\
4. ESTRATÉGIA: proponha curso de ação com justificativa\\
5. RESSALVAS: indique lacunas de informação [...]\\
\\
\#\# CALIBRAÇÃO EPISTÊMICA\\
- Use ``provavelmente'' para probabilidade > 70\%\\
- Use ``possivelmente'' para probabilidade 40--70\%\\
[...]
\end{quote}

\section{Judge Rubric}
\label{app:rubric}

The judge was prompted with the following criteria (0–8 each):

\begin{enumerate}
  \item \textbf{Reasoning depth}: Is the reasoning multi-step, explicit, and
    non-superficial?
  \item \textbf{Domain accuracy}: Are concepts, terminology, and references
    correct for the domain?
  \item \textbf{Risk identification}: Were risks identified with specificity?
  \item \textbf{Actionability}: Are recommendations specific enough to execute?
  \item \textbf{Epistemic calibration}: Does the output distinguish certainties
    from uncertainties?
\end{enumerate}

The judge received outputs with condition labels stripped (blind evaluation).

\section{Raw Results}
\label{app:raw}

Full results available as JSON in the repository. Sample from the Legal domain
(run 1):

\begin{table}[h]
\centering
\caption{Legal domain, run 1 — scores by dimension}
\begin{tabular}{lccccc|c}
\toprule
Condition & Depth & Accuracy & Risk & Action & Calib. & Total \\
\midrule
fi\_extended    & 8 & 7 & 8 & 8 & 8 & 39 \\
fi\_standard    & 7 & 7 & 7 & 7 & 7 & 35 \\
no\_fi\_extended & 7 & 7 & 7 & 7 & 7 & 35 \\
\bottomrule
\end{tabular}
\end{table}

\end{document}
